\chapter{The Thorin measure at finite resolution}\label{chap:measure}

Chapter~\ref{chap:evaluation} evaluates $\xi$ and $\xi'$ off the critical line from the P\'olya kernel alone, and so computes the Thorin density $\rho_\varepsilon$ of Section~\ref{ch:clock}, the Poisson smoothing of the zeros, at every resolution $\varepsilon>0$. The Thorin condition is RH (Proposition~\ref{prop:rigid}), and by Lagarias's criterion RH is the positivity of $\rho_\varepsilon$ at every resolution (Remark~\ref{rem:ch2:oneline}); the question is what finite computations of $\rho_\varepsilon$ can establish about it. The answer, which Chapters~\ref{chap:hcm} and \ref{chap:heat} complete, has two sides. On every window where the zeros are verified, the computed Thorin measure $\mu_\varepsilon$ converges to $U_{\RH}$, the discrete measure demanded by RH, at explicit rates as $\varepsilon\downarrow0$. Above the computed range every finite computation of the kinds in this book remains consistent with a failure of the Thorin condition, and the step that remains, positivity at every resolution and every height, is not a computational one (Section~\ref{ch:finite}).

Section~\ref{ch:motivation} asks what computation can reach and what the sine squared lemma adds, and sets the limits within which the later computations work. On the window below height $200$, where the zeros have been counted, Section~\ref{ch:distances} measures the distance between $\mu_\varepsilon$ and $U_{\RH}$ in several metrics and explains why total variation must be used at matched resolution, and Section~\ref{ch:offline} computes the signatures of an off line zero and calibrates the comparison against them. Section~\ref{ch:turing} turns the argument principle into an identity that reads the zero count and the zero density, strip by strip, from the Thorin mass. Section~\ref{ch:sin2more} develops further readings of the sine squared lemma, including Selberg's correspondence between resolution and prime cutoff near the critical line and Weil's criterion as a domination on variograms.
